% Ad Atticum 13.26 --- headnote
% ad-atticum-13-26, 14 May 45 BC, Tusculanum

\textit{Cicero to Atticus, written at the Tusculan villa on 14 May
45 BC --- Perseus dateline \textit{Scr. in Tusculano prid. Id. Mai.
a. 709 (45)}. This is the earliest letter in Book 13's mid-May to
early-June cluster, and its position late in the book is purely an
accident of ancient shelving: book 13 as transmitted opens mid-May
and the editor has placed several of the very earliest letters of
that run --- this one among them --- toward the end of the book
rather than the beginning. The grief is the recent grief of
Tullia's death; the philosophical work of the season is the
\textit{Academica} revision and the dedication-shuffle around
Varro.}

\textit{Two practical strands: a property purchase --- the
``shrine'' (\textit{fanum}) for Tullia, never named explicitly here
but unmistakable in the phrase \textit{in eius rei cupiditate quam
nosti} (``my desire for the thing you know of'') --- with
Vergilius's plot the preferred site, Clodia's second, and Drusus's
to be attempted only in desperation; and the letter Atticus had
pressed Cicero to write to Caesar, now finished, with Cicero
having second thoughts about whether it is necessary to send it.
No Greek phrases in this letter. The reference to Astura at the
opening of section 2 is to where Cicero would be, not where he is:
he is writing from Tusculum but considering whether to retreat
again to the coast, and decides instead to come on toward Rome
via Lanuvium.}
